\section{Discussion}
\label{sec:discussion}

The central empirical result is that a validated treatment layer applied as a hard retrieval gate changes \plre{} behavior under matched conditions: stale citations fall from 12.7\% to 4.1\% and F1 rises by 0.06 relative to the identical pipeline without gating (Table~\ref{tab:ablation}).
Independent checks place negative-treatment micro-F1 at 0.80 (gold spans) and 0.81 (commercial citator).
These figures are point estimates; we treat them as evidence for the mechanism rather than as significance-tested superiority.
Instance-level predictions are retained and released so that paired bootstrap CIs can be attached by any reader.

Two controls bound the claim.
First, without the gate the agentic loop alone does not improve \plre{} (0.72 vs.\ RAG-only 0.73), so gains are not explained by extra test-time compute.
Second, a surface-form heuristic recovers many explicit \textsc{Overrules}/\textsc{Supersedes} mentions but almost no material narrowing, and citator-oracle intervals flip only 3.8\% of held-out pairs---the extracted layer already tracks the commercial reference for downstream decisions.
We deliberately do not evaluate pure date filters or learned temporal re-rankers on \plre{} in the present study; a date filter cannot recover supersession endpoints from later opinions, which is the extraction problem this work targets.
External-stack transfer (Protocol~P1) remains a pilot and is not claimed.

On severity triage ($n = 82$), \sysname{} leads iteration-matched B5 by a small margin (0.79 vs.\ 0.75); that comparison is secondary to \plre{}.
